% The spoken body of the homily for the Seventeenth Sunday after Pentecost.
% Continuous preaching, ready to read aloud: no heading, no direction to the
% preacher, no citation inside the speech. The loci, the relation to the
% reviewed readings, the audience, the word count and the pace estimate stand
% in the terminal note, which no one reads out. The blank lines are the
% paragraph breaks a speaker recovers his place by, and the four vertical
% spaces divide the speech into its five movements: the three questions and
% the silence; the double commandment; Christ's own question; the Epistle's
% unity and the Collect that asks for what the Gospel commands; and the
% Introit, the Offertory, the Creed, the Communion and the Postcommunion.

Three questions are asked in today's Gospel, and the last of them is never
answered. A doctor of the law asks the first, to test him: ``Master, which is
the great commandment in the law?'' He is given an answer, and then a second
answer he had not asked for. Our Lord asks the second himself: ``What think
you of Christ? Whose son is he?'' That one the men who had come to trap him
can answer, and they answer it rightly: ``David's.'' So he asks the third.
David, he says, calls the Christ his Lord, and says it in the Spirit: ``If
David then call him Lord, how is he his son?'' Saint Matthew closes the scene
with a sentence almost painful to read aloud: ``And no man was able to answer
him a word: neither durst any man from that day forth ask him any more
questions.''

Hold on to that silence. Everything else in this Mass is standing around it.

\bigskip

The first answer we think we already know by heart. ``Thou shalt love the
Lord thy God with thy whole heart and with thy whole soul and with thy whole
mind.'' Three times, whole. Saint Augustine asks what that word is doing
there three times over, and his answer is plain and uncomfortable. It means,
he says, ``that no part of our life is to be unoccupied, and to afford room,
as it were, for the wish to enjoy some other object.'' He is not saying that
we may love nothing else. He is saying that everything else has to be loved
inside that love and carried along by it, the way one current carries
whatever falls into it. There is no second channel. Whatever we insist on
keeping out of that one current we are not really loving at all; we have only
made a small god of it.

And then the second commandment, which nobody had asked for: ``Thou shalt
love thy neighbour as thyself.'' Why is it ``like to this''? Saint John
Chrysostom gives the reason in a line: ``this makes the way for that, and by
it is again established.'' The love of the neighbour opens the road to the
love of God; and the love of God, in its turn, re-establishes the love of the
neighbour. Each holds the other up. The first commandment cannot be
obeyed alone, in private, with the door shut.

\bigskip

Now the second question, which looks at first like a change of subject. We
were speaking about love, and suddenly we are inside an old psalm about a
throne and a footstool. Chrysostom reads the two halves of this Gospel as one
lesson. When Scripture says that there is one God, he writes, it is said
``not to the rejection of the Son, but to make the distinction from idols'';
and so our Lord puts the question that will not let anyone keep the
commandment and lose him. He was leading them, Chrysostom says, ``to confess
Him also to be God.''

Their answer was true and one step short. Saint Augustine,
preaching on the very psalm our Lord quotes back at them, tells his own people
that they must answer as those men answered --- the Christ is David's son ---
and then he tells them not to remain where those men remained. In the form of
God, he is David's Lord. Taking the form of a servant, he is David's son.
That is the answer nobody in the Temple would give. And there is the reason
the commandment and the question stand together in one Gospel. First we are
told whom we must love with everything we have. Then we are shown who he is. The God of the great
commandment has a face, and it is the face of the man asking the question.

\bigskip

Saint Paul, who writes as a prisoner, tells us what that love looks like
when a great many people attempt it at once. It is not a feeling. It is a
manner of life: ``With all humility and mildness, with patience, supporting
one another in charity. Careful to keep the unity of the Spirit in the bond
of peace.'' Notice the verb. He does not tell the Ephesians to make their
unity. He tells them to keep it, because it has already been given to them:
``One body and one Spirit: as you are called in one hope of your calling. One
Lord, one faith, one baptism. One God and Father of all.'' Seven times in
three verses he says the word one, to a church he is urging not to come
apart. The one God whom we are to love with an undivided heart has one body in
this world, and we are looking at part of it.

``Supporting one another in charity'' is not goodwill toward humanity in
general. In every parish it has a name and a face: the person who talks too
long after Mass, or the one who never talks at all; the one who was unfair to
you and has never noticed; the one whose grief, or whose opinions, you can
bear for about four minutes. That person is where the first
commandment is actually being kept, or actually being refused. There is no
other place for it to happen this week.

Which is why the Church, at the beginning of this Mass, does not command it
of us. She asks for it: ``Grant, we beseech thee, O Lord, that thy people may
avoid all contagion of the devil: and with a clean heart follow thee, the
only true God.'' Kept clear of what corrupts us, and following God alone:
that is the one current again, and no second channel. The prayer asks of God
what the Gospel commands of us. The commandment comes to us as a command and
goes back to God as a petition,
because we cannot manufacture it ourselves.

\bigskip

And notice the tone in which this Mass asks. At the door it does not ask for
justice. ``Thou art just, O Lord: and thy judgment is right'' --- that much
is confessed at once --- and then: ``Deal with thy servant according to thy
mercy.'' According to thy mercy: not according to mine. Later, when the gifts
are carried up, the Church borrows the words of Daniel, praying in exile for
a desolate sanctuary and for the people that carries God's name. Saint Jerome
noticed that Daniel confesses that people's sins in the first person, and
gave the reason in five words: \latin{quia unus e populo est}, because he is
one of the people. That is the posture of everybody in this church today: not
a private soul with a private account to settle, but one of a people, praying
for the whole of it and included in everything it confesses.

Which brings us back to the silence. ``No man was able to answer him a
word.'' The question behind that silence was never rhetorical, and it did not
stop being asked when the Pharisees walked away. How can the one David calls
Lord be David's son? In a few moments we are going to do something about that
silence. We will stand and say the Creed, which is the answer those men would
not give; then bread and wine will be carried to this altar and offered, and
we shall offer with them, and receive the one we have just confessed.

At Communion the Church will sing an old verse over us: ``Vow ye, and pay to
the Lord your God: all you that are round about him bring presents.'' Vow,
and pay. That is a good rule for the coming week, provided we do not vow
much. Vow one thing that can actually be paid. One person you will carry
instead of avoid. One hour that has been quietly taking up the room in your
heart that belongs to God, which you will give back to him. Then pay it.

The prayer after Communion asks that ``our vices be cured'' and that ``an
eternal remedy'' be given us. Eternal: what is begun at this altar is not
finished at this altar. It is finished in the one hope of our calling: where
the whole heart is at last whole, and where the question our Lord asked in
the Temple is answered by sight and no longer by faith. Until that day we answer it as the Church
answers it --- standing, in the Creed; kneeling, when he is given to us; and
during the week, in the one person we would rather not carry.
